% ==========================================================================
% DELIVERABLE 2: COMPLETE PAPER OUTLINE (Sections 1-8)
%
% Title: Isomorph-Eval: Separating Reasoning from Recall in LLM Evaluation
%        via Structurally Equivalent Benchmarks and Item Response Theory
%
% Target: NeurIPS 2027 Evaluations & Datasets Track
% ==========================================================================
%
% OUTLINE SUMMARY
%
% Section 1: Introduction (1.5 pages) — FULL TEXT BELOW
%
% Section 2: Related Work (1.5 pages)
%   2.1 Benchmark Contamination Detection
%       GSM1K [Zhang NeurIPS D&B 2024]: parallel benchmark, 8-13% drops.
%       ConStat [Dekoninck NeurIPS 2024]: statistical test + reference models.
%       MMLU-CF [2024]: closed contamination-free MMLU variant.
%       PaCoST [2024]: paired confidence testing (needs model internals).
%   2.2 Living / Dynamic Benchmarks
%       LiveBench [White ICLR 2025 Spotlight]: monthly-updated questions.
%       FrontierMath [Glazer 2024]: expert-curated hard math.
%       EternalMath [2026]: living benchmark from new theorems.
%       Key gap: all create NEW benchmarks, none generate isomorphs of EXISTING ones.
%   2.3 IRT for LLM Evaluation
%       Rodriguez [ACL 2021], Polo [ICML 2024], Zhou PSN-IRT [AAAI 2026 Oral],
%       Singhal MEDIRT [2025], Uebayashi M3IRT [2026], LaRT [2025].
%       Key gap: all fix aggregation, none address contamination.
%   2.4 Evaluation Failure Scaling Law
%       Our prior work [Kang 2026]: S x D interaction, gamma_3 = +0.199.
%       This paper adds the third axis: contamination C.
%   Feature comparison table: 7 dimensions x 5 methods.
%
% Section 3: Background and Notation (0.75 pages)
%   3.1 2PL Item Response Theory (model, estimation, missing data)
%   3.2 Differential Item Functioning (DIF) — the psychometric bridge
%   3.3 EFSL recap: 1 - rho = gamma_0 + gamma_1*S + gamma_2*D + gamma_3*S*D
%
% Section 4: The Isomorph-Eval Framework (2.5 pages)
%   4.1 Task Isomorphism (Definitions 1-2, Conjecture 1)
%   4.2 The Isomorphic Engine (Parser -> Mutator -> Generator -> Verifier)
%   4.3 The Contamination Delta (Definitions 3-6, hypothesis test)
%   4.4 Diagnostic Archetypes (4 types, threshold table)
%   4.5 The Three-Body Problem: 1 - rho = f(S, D, C)
%
% Section 5: Experimental Setup (1 page)
%   5.1 Dataset: GSM8K (667 problems x 5 variants = 3,866 items; 100-item stratified subset)
%   5.2 Models: Llama-3.1-8B, Llama-4-Scout-17B, Qwen3-32B
%   5.3 Protocol: T=1 trial, temp=0.0, CoT prompt, boxed extraction, bootstrap CIs
%   5.4 Baselines: raw accuracy gap (GSM1K), our IRT-weighted
%
% Section 6: Results (2 pages)
%   6.1 Main Results — Table + Figure 1 (diagnostic bar chart)
%   6.2 Three-Body Surface — Figure 2 (4-panel contour/bar)
%   6.3 Per-Item DIF Analysis — which items are memorized?
%   6.4 Reasoning Structure Hypothesis validation — b_orig vs b_iso correlation
%
% Section 7: Discussion (1 page)
%   7.1 Limitations (parser coverage, surface effects, domain scope)
%   7.2 The Goodhart Problem (Archetype 4: Latent Contaminator)
%   7.3 Implications (leaderboards, regulation, infrastructure)
%   7.4 Connection to EFSL (the unified theory)
%
% Section 8: Conclusion (0.5 pages)
%   Summary, open-source release, future work (logic, code, Physical AI)
%
% ==========================================================================


% ==========================================================================
% DELIVERABLE 3: FULL INTRODUCTION TEXT (Section 1)
%
% Ready to copy-paste into any NeurIPS LaTeX template.
% ==========================================================================

\section{Introduction}

The credibility of AI evaluation depends on two premises that are
increasingly difficult to defend. The first is that benchmark items have
not leaked into training data. The second is that aggregating scores by
simple averaging produces rankings that faithfully reflect true system
capability. A growing body of evidence suggests that both premises are
violated in practice, and that the violations interact in ways that are
worse than either alone.

\paragraph{The contamination crisis.}
Test set contamination---the inadvertent or deliberate inclusion of
benchmark items in training corpora---has emerged as a fundamental threat
to evaluation integrity. Zhang et al.~\cite{zhang2024gsm1k} constructed
GSM1K, a parallel benchmark to GSM8K, and found that several prominent
open-source models exhibited performance drops of 8--13\%, with some
models showing near-perfect accuracy on original items while failing on
structurally similar but previously unseen variants. Dekoninck et
al.~\cite{dekoninck2024constat} proposed ConStat, a statistical test
that flags contamination via performance gaps on rephrased items, using
reference models as calibration anchors. Concurrently, MMLU-CF
\cite{mmlucf2024} demonstrated that GPT-4o's score drops from 88.0\% on
the original MMLU to 73.4\% on a contamination-free variant---a 14.6
percentage-point gap that cannot be attributed to task difficulty alone.
These findings indicate that a non-trivial fraction of the performance
gains reported on widely-used benchmarks may reflect memorization rather
than generalization.

\paragraph{The aggregation crisis.}
Independently, the methodology by which benchmark scores are aggregated
has received increasing scrutiny. In prior work, we introduced the
\emph{Evaluation Failure Scaling Law} (EFSL)~\cite{kang2026efsl},
demonstrating through controlled cross-domain simulations that Spearman
rank correlation between simple-average rankings and ground-truth
rankings degrades from $\rho = 1.000$ to $\rho = 0.242$ as the product
of data sparsity $S$ and item difficulty gap $D$ increases under biased
missingness. A 150-condition grid sweep confirmed a strong and
statistically significant $S \times D$ interaction
($\gamma_3 = +0.199$, $t = 13.05$), while two-parameter logistic (2PL)
Item Response Theory (IRT) maintained $\rho \geq 0.993$ across all
conditions. Concurrently, Zhou et al.~\cite{zhou2026psnirt} conducted
the most comprehensive IRT analysis of LLM benchmarks to date, revealing
significant quality shortcomings across 11 benchmarks comprising 41,871
items; Singhal et al.~\cite{singhal2025medirt} showed that IRT-based
rankings outperform accuracy-based rankings across six external medical
benchmarks; and recent work on latency-aware IRT~\cite{lart2025}
extended the framework to jointly model response time and accuracy.

\paragraph{The compound failure.}
While each crisis has received individual attention, no prior work has
examined their joint interaction. We identify what we term the
\emph{Three-Body Problem of AI evaluation}: the compound failure that
arises when data sparsity, item difficulty heterogeneity, and benchmark
contamination co-occur. Consider a concrete scenario: an open-source
model has memorized 15\% of GSM8K through training data leakage and is
evaluated on a sparse subset of items (60\% coverage) that
disproportionately includes the easy items it has memorized. Its simple
average score is inflated by all three mechanisms simultaneously: (1)
sparsity prevents the evaluation from sampling its failure modes; (2)
difficulty bias means the tested items are disproportionately easy; (3)
contamination means even the correctly-tested items receive a
memorization boost. The resulting ranking is not merely noisy---it is
\emph{systematically} biased in a direction that rewards contaminated
models evaluated under favorable conditions.

IRT alone corrects for (1) and (2) but not (3), because it cannot
distinguish genuine ability from contamination-inflated performance on
the same items. Contamination detection alone (GSM1K, ConStat) addresses
(3) but provides no correction for (1) and (2), and its difficulty
adjustment relies on generic regression rather than the formal item
parameter structure that IRT provides. Neither approach alone solves the
compound problem.

\paragraph{This paper: Isomorph-Eval.}
We present \textbf{Isomorph-Eval}, a unified framework that addresses
all three failure axes simultaneously. The framework consists of three
components:

\begin{enumerate}[leftmargin=*, nosep]
\item \textbf{$\boldsymbol{\tau}$-isomorphism} (Section~4.1): a
graph-theoretic definition of structural equivalence between benchmark
items. Two items are $\tau$-isomorphic if there exists a graph
isomorphism between their reasoning graphs that preserves operation
types, dependency structure, and complexity weights, while ensuring
complete surface novelty. Unlike prior rephrasing approaches that rely
on semantic similarity heuristics, $\tau$-isomorphism is verified
through a deterministic structural comparison, ensuring that generated
variants test identical cognitive operations with known-correct answers.

\item \textbf{The Isomorphic Engine} (Section~4.2): an open-source
pipeline that automatically generates unlimited verified
$\tau$-isomorphic variants of any binary-scored benchmark item. The
engine decomposes items into reasoning graphs via a hybrid rule-based
and LLM-structured parser (achieving 57\% coverage on GSM8K with
zero LLM calls, and approximately 88\% with structured extraction),
mutates entity and value bindings while locking structural invariants,
and verifies output fidelity through round-trip re-parsing. A key
architectural contribution is the \emph{template-first surface
generator}, in which the LLM functions exclusively as a surface
realizer---never as a reasoner---with its output validated by forward
graph execution.

\item \textbf{The Contamination Delta}
$\boldsymbol{\Delta_{\mathrm{contam}}^{\mathrm{IRT}}}$ (Section~4.3):
an IRT-grounded diagnostic metric that quantifies memorization by
comparing model performance on original items versus isomorphic
variants, weighted by item discrimination parameters estimated from
clean (isomorphic) data. Grounded in the psychometric framework of
Differential Item Functioning (DIF)~\cite{lord1980,holland1993}, the
metric provides per-item Wald tests with Bonferroni-corrected confidence
intervals and a global likelihood ratio test, classifying models into
four diagnostic archetypes: \emph{Robust Reasoner}, \emph{Pure
Memorizer}, \emph{Syntactic Matcher}, and \emph{Latent Contaminator}.
In controlled simulations, the IRT-weighted delta is 23$\times$ more
discriminative than the raw accuracy gap used by prior methods.
Empirical evaluation on five LLMs (including a negative control with
no GSM8K in its training data) on 59 verified-clean GSM8K items shows
that four of five models exhibit no significant contamination
($\Delta_{\mathrm{raw}} \in [+0.049, +0.066]$, all $p > 0.05$),
while Llama~3.1 8B shows a mild gap ($\Delta_{\mathrm{raw}} = +0.126$,
$p = 0.002$). Critically, the framework also uncovered a data quality
confound in annotation-chain replay that inflated apparent
contamination signals uniformly across all models, demonstrating the
importance of independent answer verification in benchmark construction.
\end{enumerate}

\paragraph{The unified theory.}
Isomorph-Eval extends our prior work on the Evaluation Failure Scaling
Law~\cite{kang2026efsl} by adding contamination as a third axis of the
failure surface. Where EFSL established that
$1 - \rho_{\mathrm{avg}} = f(S, D)$ with a strong $S \times D$
interaction, we show that the full model is
$1 - \rho = f(S, D, C)$, where $C$ measures differential contamination
across models. The EFSL provides the correction mechanism for $S$ and
$D$ (2PL IRT estimation); Isomorph-Eval provides the correction
mechanism for $C$ (isomorphic items with verified structural
equivalence). Together, they constitute what we believe to be the first
unified framework for fair AI evaluation under the joint presence of
data sparsity, difficulty heterogeneity, and benchmark contamination.

\paragraph{Contributions.}
Our contributions are as follows:

\begin{enumerate}[leftmargin=*, nosep]
\item We formalize $\tau$-isomorphism as a verifiable graph-theoretic
condition for benchmark item equivalence, grounding the concept of
``structurally identical problem'' in operation-type and
complexity-weight preservation over reasoning graphs.

\item We design and release the Isomorphic Engine, a four-stage pipeline
(Parser $\to$ Mutator $\to$ Generator $\to$ Verifier) that produces
unlimited novel benchmark variants with automatically computed answers
and deterministic structural verification.

\item We define
$\Delta_{\mathrm{contam}}^{\mathrm{IRT}}$, an IRT-weighted
contamination metric grounded in the DIF framework, and demonstrate its
superior discriminative power over raw accuracy gaps.

\item We identify the Three-Body Problem of AI evaluation and show that
the compound interaction of sparsity, difficulty heterogeneity, and
contamination produces ranking failures that neither IRT alone nor
contamination detection alone can correct---but that the unified
framework recovers.
\end{enumerate}

\noindent We release the complete pipeline---parser, mutator, generator,
verifier, contamination metric, and evaluation runner---as open-source
software requiring only NumPy, SciPy, and a standard Python
environment.\footnote{Code and data:
\url{https://github.com/testofschool/isomorph-eval}}
